%============================================================
%  DISPENSA: Computational Mathematics for Learning and Data Analysis
%  Università di Pisa — Laurea Magistrale in Informatica
%  Prof. Antonio Frangioni (Ottimizzazione) &
%  Prof. Alice Cortinovis (Analisi Numerica Lineare)
%  A.A. 2025/2026
%============================================================
\documentclass[11pt, a4paper, twoside]{book}

%--- Encoding & Language ---
\usepackage[T1]{fontenc}
\usepackage[utf8]{inputenc}
\renewcommand{\contentsname}{Indice}
\renewcommand{\chaptername}{Capitolo}
\renewcommand{\appendixname}{Appendice}
\renewcommand{\bibname}{Bibliografia}
\renewcommand{\figurename}{Figura}
\renewcommand{\tablename}{Tabella}
\renewcommand{\partname}{Parte}

%--- Math ---
\usepackage{amsmath, amssymb, amsthm, mathtools}
\usepackage{bm}

%--- Layout ---
\usepackage[top=2.5cm, bottom=2.5cm, left=3cm, right=2.5cm]{geometry}
\usepackage{microtype}
\usepackage{setspace}
\onehalfspacing

%--- Fonts ---
\usepackage{lmodern}

%--- Colors & Graphics ---
\usepackage[dvipsnames]{xcolor}
\usepackage{graphicx}
\usepackage{tikz}
\usepackage{pgfplots}
\pgfplotsset{compat=1.18}

%--- Hyperlinks ---
\usepackage[colorlinks=true, linkcolor=NavyBlue, citecolor=Green, urlcolor=RoyalBlue]{hyperref}

%--- Code listings ---
\usepackage{listings}
\lstset{
  basicstyle=\ttfamily\small,
  keywordstyle=\color{NavyBlue}\bfseries,
  commentstyle=\color{gray},
  stringstyle=\color{OliveGreen},
  numbers=left, numberstyle=\tiny\color{gray},
  breaklines=true,
  frame=single, framerule=0.4pt,
  backgroundcolor=\color{gray!8},
  tabsize=2
}

%--- Fancy headers ---
\usepackage{fancyhdr}
\setlength{\headheight}{14pt}
\pagestyle{fancy}
\fancyhf{}
\fancyhead[LE]{\small\textit{\leftmark}}
\fancyhead[RO]{\small\textit{\rightmark}}
\fancyfoot[C]{\thepage}
\renewcommand{\headrulewidth}{0.4pt}

%--- TiKZ styles ---
\usetikzlibrary{arrows.meta, positioning, shapes.geometric, calc}

%--- Table of contents depth ---
\setcounter{tocdepth}{2}
\setcounter{secnumdepth}{3}

%=== Custom Theorem Environments ===
\theoremstyle{definition}
\newtheorem{defn}{Definizione}[chapter]
\newtheorem{example}{Esempio}[chapter]
\newtheorem{remark}{Osservazione}[chapter]

\theoremstyle{plain}
\newtheorem{thm}{Teorema}[chapter]
\newtheorem{prop}{Proposizione}[chapter]
\newtheorem{lemma}{Lemma}[chapter]
\newtheorem{corol}{Corollario}[chapter]

\theoremstyle{remark}
\newtheorem{nota}{Nota}[chapter]
\newtheorem{intuition}{Intuizione}[chapter]

%=== Custom boxes ===
\usepackage{tcolorbox}
\tcbuselibrary{theorems, skins, breakable}

\newtcolorbox{keybox}[1][]{
  enhanced, breakable,
  colback=NavyBlue!5, colframe=NavyBlue!60,
  fonttitle=\bfseries,
  title={#1}
}
\newtcolorbox{algobox}[1][]{
  enhanced, breakable,
  colback=OliveGreen!5, colframe=OliveGreen!50,
  fonttitle=\bfseries,
  title={Algoritmo: #1}
}
\newtcolorbox{warningbox}{
  enhanced,
  colback=Bittersweet!8, colframe=Bittersweet!70,
  fonttitle=\bfseries,
  title={Attenzione}
}

%=== Math operators ===
\DeclareMathOperator*{\argmin}{arg\,min}
\DeclareMathOperator*{\argmax}{arg\,max}
\DeclareMathOperator{\rank}{rank}
\DeclareMathOperator{\dom}{dom}
\DeclareMathOperator{\conv}{conv}
\DeclareMathOperator{\diag}{diag}
\DeclareMathOperator{\tr}{tr}
\DeclareMathOperator{\sign}{sign}
\DeclareMathOperator{\cond}{cond}
\DeclareMathOperator{\prox}{prox}
\newcommand{\norm}[1]{\left\lVert #1 \right\rVert}
\newcommand{\abs}[1]{\left\lvert #1 \right\rvert}
\newcommand{\R}{\mathbb{R}}
\newcommand{\N}{\mathbb{N}}
\newcommand{\grad}{\nabla}
\newcommand{\eps}{\varepsilon}

%============================================================
\begin{document}

%=== TITLE PAGE ===
\begin{titlepage}
\centering
\vspace*{1.5cm}
{\LARGE \textbf{Università di Pisa}\\[0.3em]
\large Laurea Magistrale in Informatica}

\vspace{1cm}
\rule{\textwidth}{1.5pt}\\[0.4em]
{\Huge \textbf{Computational Mathematics}\\[0.4em]
\huge \textbf{for Learning and Data Analysis}}\\[0.2em]
\rule{\textwidth}{1.5pt}

\vspace{0.8cm}
{\Large Dispensa del corso}

\vspace{0.5cm}
{\large
\textit{Basata sulle lezioni di}\\[0.3em]
\textbf{Prof. Antonio Frangioni} (Ottimizzazione)\\
\textbf{Prof. Alice Cortinovis} (Analisi Numerica Lineare)\\[0.5em]
A.A.\ 2025/2026
}

\vfill
\begin{tcolorbox}[colback=NavyBlue!10, colframe=NavyBlue!50, width=0.85\textwidth]
\centering\small
Questa dispensa integra le slide ufficiali, le note di Frangioni (2025/26)
e gli appunti dalle esercitazioni di Cortinovis.
È redatta in italiano con terminologia tecnica in inglese.\\
\textit{Obiettivo: copertura completa di tutto il programma del corso.}
\end{tcolorbox}
\vspace{0.5cm}
{\small Versione 2.0 — Aprile 2026}
\end{titlepage}

\tableofcontents
\clearpage

%============================================================
\chapter*{Introduzione al Corso}
\addcontentsline{toc}{chapter}{Introduzione al Corso}

Il corso di \textbf{Computational Mathematics for Learning and Data Analysis} è tenuto da due professori:
\begin{itemize}
  \item \textbf{Prof. Antonio Frangioni} — parte di \emph{ottimizzazione}: risoluzione di $\min_{x \in S} f(x)$ per classi crescenti di funzioni e insiemi ammissibili.
  \item \textbf{Prof. Alice Cortinovis} — parte di \emph{analisi numerica lineare (NLA)}: risoluzione efficiente e accurata del problema dei minimi quadrati $\min_{x} \|Ax - b\|^2$ e dei sistemi lineari.
\end{itemize}

Le due parti sono strettamente interconnesse: molti algoritmi di ottimizzazione richiedono la risoluzione di sistemi lineari come sotto-passo (es.\ il passo di Newton). Capire entrambe è fondamentale.

\medskip
\textbf{Prerequisiti:} algebra lineare (spazi vettoriali, autovalori), analisi (Taylor, derivate parziali), nozioni di programmazione (Matlab/Python).

%============================================================
\part{Ottimizzazione Matematica}

%------------------------------------------------------------
\chapter{Semplici Problemi di Ottimizzazione}
%------------------------------------------------------------

\section{Formulazione generale}

\begin{defn}[Problema di ottimizzazione]
Dato un insieme \emph{ammissibile} $S \subseteq \R^n$ e una funzione \emph{obiettivo} $f: S \to \R$, il problema è:
\[
  f^* = \min_{x \in S} f(x).
\]
Un punto $x^* \in S$ tale che $f(x^*) = f^*$ è una \emph{soluzione ottima} (o \emph{minimizzatore}).
\end{defn}

\begin{remark}
Il valore ottimo $f^*$ è unico se esiste, ma possono esistere più minimizzatori $x^*$. Massimizzare $g(x)$ equivale a minimizzare $f(x) = -g(x)$. Un punto $x_\eps$ è una \emph{soluzione $\eps$-ottima} se $f(x_\eps) \leq f^* + \eps$.
\end{remark}

\subsection{Minimi locali vs globali}

\begin{defn}
$x^*$ è un \emph{minimo locale} se esiste $\delta > 0$ tale che $f(x^*) \leq f(x)$ per ogni $x \in S$ con $\|x - x^*\| \leq \delta$. È un \emph{minimo globale} se $f(x^*) \leq f(x)$ per ogni $x \in S$.
\end{defn}

Per funzioni convesse, ogni minimo locale è globale. Per funzioni non convesse, possono esistere minimi locali che non sono globali — trovare il globale è in generale NP-difficile.

%------------------------------------------------------------
\section{Ottimizzazione univariata (1D)}

Il caso $n=1$ è fondamentale perché molte line search si riducono a questo.

\begin{defn}[Funzione unimodale]
$f: [a,b] \to \R$ è unimodale se esiste $x^* \in [a,b]$ tale che $f$ è strettamente decrescente su $[a, x^*]$ e strettamente crescente su $[x^*, b]$.
\end{defn}

\subsection{Golden Section Search}

\begin{algobox}[Golden Section Search]
Dato $f$ unimodale su $[a, b]$, con $\tau = (\sqrt{5}-1)/2 \approx 0.618$:
\begin{enumerate}
  \item $x_1 = a + (1-\tau)(b-a)$, $x_2 = a + \tau(b-a)$.
  \item Se $f(x_1) < f(x_2)$: poni $b \leftarrow x_2$, altrimenti $a \leftarrow x_1$.
  \item Aggiorna il punto interno rimanente senza ricalcolarlo (1 valutazione per passo).
  \item Ripeti finché $|b-a| \leq \eps$.
\end{enumerate}
Dopo $k$ passi, l'intervallo si è ridotto a $(b-a) \cdot \tau^k$. Costo: $O(\log(1/\eps))$ valutazioni.
\end{algobox}

\begin{intuition}
La scelta di $\tau = (\sqrt{5}-1)/2$ (rapporto aureo) garantisce che il punto ``risparmiato'' cada nella posizione giusta nel sotto-intervallo successivo: un'elegante proprietà di auto-similarità.
\end{intuition}

\subsection{Metodo di Newton 1D}

Per $f \in C^2$, partendo da $x_0$:
\[
  x_{k+1} = x_k - \frac{f'(x_k)}{f''(x_k)}.
\]
Convergenza \emph{quadratica} ($\|e_{k+1}\| = O(\|e_k\|^2)$) vicino a una radice di $f'$ con $f''(x^*) \neq 0$.

%------------------------------------------------------------
\chapter{Ottimizzazione Non Vincolata}
%------------------------------------------------------------

\section{Richiami di analisi in più variabili}

\subsection{Gradiente, Hessiana, sviluppo di Taylor}

\begin{defn}[Gradiente]
Per $f: \R^n \to \R$ differenziabile:
\[
  \nabla f(x) = \left(\frac{\partial f}{\partial x_1}, \ldots, \frac{\partial f}{\partial x_n}\right)^\top \in \R^n.
\]
$\nabla f(x)$ punta nella \emph{direzione di massima crescita} di $f$ in $x$.
\end{defn}

\begin{defn}[Hessiana]
Per $f \in C^2(\R^n)$:
\[
  \nabla^2 f(x) = H_f(x) \in \R^{n \times n}, \quad [H_f(x)]_{ij} = \frac{\partial^2 f}{\partial x_i \partial x_j}(x).
\]
Per il teorema di Schwarz, $H_f$ è simmetrica.
\end{defn}

\textbf{Sviluppo di Taylor al secondo ordine:}
\[
  f(x+d) = f(x) + \nabla f(x)^\top d + \tfrac{1}{2} d^\top H_f(x) d + o(\|d\|^2).
\]

\begin{example}[Funzione quadratica]
$f(x) = \frac{1}{2}x^\top Q x - c^\top x$ con $Q$ simmetrica: $\nabla f(x) = Qx - c$, $H_f(x) = Q$ (costante).
\end{example}

\subsection{Condizioni di ottimalità}

\begin{thm}[Condizioni del primo e secondo ordine]
Per $f \in C^2$:
\begin{enumerate}
  \item \textbf{Necessaria 1° ordine}: Se $x^*$ è minimo locale, $\nabla f(x^*) = 0$ (punto stazionario).
  \item \textbf{Necessaria 2° ordine}: Se $x^*$ è minimo locale, $H_f(x^*) \succeq 0$.
  \item \textbf{Sufficiente}: Se $\nabla f(x^*) = 0$ e $H_f(x^*) \succ 0$, allora $x^*$ è minimo locale isolato.
\end{enumerate}
\end{thm}

\subsection{Classi di regolarità importanti}

\begin{defn}[$L$-smooth]
$f \in C^1$ è \emph{$L$-smooth} se $\nabla f$ è Lipschitz con costante $L > 0$:
\[
  \|\nabla f(x) - \nabla f(y)\| \leq L\|x - y\| \quad \forall x, y.
\]
Equivalentemente (per $f \in C^2$): $\nabla^2 f(x) \preceq LI$ per ogni $x$.
\end{defn}

\begin{prop}[Lemma di discesa]
Se $f$ è $L$-smooth, allora per ogni $x, d$:
\[
  f(x + d) \leq f(x) + \nabla f(x)^\top d + \frac{L}{2}\|d\|^2.
\]
Questo è il \emph{modello quadratico superiore} di $f$.
\end{prop}

\begin{defn}[$\tau$-strongly convex]
$f$ è \emph{fortemente convessa} con modulo $\tau > 0$ se:
\[
  f(y) \geq f(x) + \nabla f(x)^\top(y-x) + \frac{\tau}{2}\|y-x\|^2 \quad \forall x, y.
\]
Equivalentemente: $\nabla^2 f(x) \succeq \tau I$, cioè tutti gli autovalori dell'Hessiana sono $\geq \tau$.
\end{defn}

\begin{remark}
La classe ``più conveniente'' per l'ottimizzazione è $f$ $L$-smooth e $\tau$-fortemente convessa, con $\kappa = L/\tau$ (numero di condizionamento dell'ottimizzazione). Se $\kappa \gg 1$, il problema è \emph{difficile}.
\end{remark}

\subsection{Funzioni convesse}

\begin{defn}[Funzione convessa]
$f$ è convessa se per ogni $x, y$ e $\lambda \in [0,1]$:
\[
  f(\lambda x + (1-\lambda) y) \leq \lambda f(x) + (1-\lambda) f(y).
\]
\end{defn}

\begin{prop}[Caratterizzazioni equivalenti per $f \in C^2$]
Le seguenti sono equivalenti: (a) $f$ convessa; (b) $f(y) \geq f(x) + \nabla f(x)^\top(y-x)$ (disuguaglianza del gradiente); (c) $H_f(x) \succeq 0$ per ogni $x$.
\end{prop}

\begin{thm}[Ottimalità globale]
Se $f$ è convessa, ogni minimo locale è globale. Se $f$ è strettamente convessa, il minimizzatore è unico.
\end{thm}

%------------------------------------------------------------
\section{Metodi del gradiente (smooth)}

\subsection{Gradient Descent a passo fisso}

\begin{algobox}[Gradient Descent]
Dato $x_0 \in \R^n$, step size $\alpha > 0$:
\[
  x_{k+1} = x_k - \alpha \nabla f(x_k), \quad k = 0, 1, 2, \ldots
\]
\end{algobox}

\begin{intuition}
$-\nabla f(x_k)$ è la direzione di massima discesa. Muoversi di un passo $\alpha$ in tale direzione riduce $f$ di $\alpha \|\nabla f(x_k)\|^2$ (approssimativamente), se $\alpha$ è abbastanza piccolo.
\end{intuition}

\subsubsection{Convergenza per funzioni quadratiche}

\begin{thm}[Convergenza GD per $f$ quadratica $Q \succ 0$]
Sia $f(x) = \frac{1}{2}x^\top Q x - c^\top x$ con $0 < \lambda_{\min} = \lambda_n \leq \lambda_1 = \lambda_{\max}$ autovalori di $Q$. Con passo ottimo $\alpha^* = \frac{2}{\lambda_{\min} + \lambda_{\max}}$:
\[
  \|x_k - x^*\|_Q \leq \left(\frac{\kappa(Q) - 1}{\kappa(Q) + 1}\right)^k \|x_0 - x^*\|_Q,
\]
dove $\kappa(Q) = \lambda_{\max}/\lambda_{\min}$ e $\|v\|_Q^2 = v^\top Q v$.
\end{thm}

\begin{remark}[Disuguaglianza di Kantorovich]
Nell'analisi di convergenza, la chiave è la \emph{disuguaglianza di Kantorovich}:
\[
  \frac{(g^\top g)^2}{(g^\top Qg)(g^\top Q^{-1}g)} \geq \frac{4\lambda_{\min}\lambda_{\max}}{(\lambda_{\min}+\lambda_{\max})^2} = 1 - \left(\frac{\kappa-1}{\kappa+1}\right)^2.
\]
Essa misura quanto la curvatura non uniforme di $f$ rallenti la discesa.
\end{remark}

\subsubsection{Convergenza per funzioni smooth convesse}

\begin{thm}[GD con passo $1/L$: convergenza $O(1/k)$]
Sia $f$ convessa e $L$-smooth. Con $\alpha = 1/L$:
\[
  f(x_k) - f^* \leq \frac{L\|x_0 - x^*\|^2}{2k}.
\]
\end{thm}

\textit{Prova (sketch):} Dal lemma di discesa con $\alpha = 1/L$:
\[
  f(x_{k+1}) \leq f(x_k) - \frac{1}{2L}\|\nabla f(x_k)\|^2.
\]
Sommando $k$ passi e usando la disuguaglianza del gradiente, si ottiene la tesi.

\begin{thm}[GD per funzioni fortemente convesse: convergenza lineare]
Se $f$ è $L$-smooth e $\tau$-fortemente convessa, con $\alpha = 1/L$:
\[
  \|x_k - x^*\|^2 \leq \left(1 - \frac{\tau}{L}\right)^k \|x_0 - x^*\|^2 = \left(1 - \frac{1}{\kappa}\right)^k \|x_0 - x^*\|^2.
\]
La convergenza è lineare (geometrica) con fattore $\rho = 1 - 1/\kappa$.
\end{thm}

\begin{remark}[Confronto tra i tassi]
$O(1/k)$ (convex) vs $O(\rho^k)$ (strongly convex): il secondo è \emph{esponenzialmente} più veloce. Tuttavia, per $\kappa$ grande, anche la convergenza lineare è molto lenta ($\rho \approx 1 - 1/\kappa$ piccolo).
\end{remark}

\subsubsection{Metodi con momento: Nesterov}

Il metodo di Nesterov (1983) aggiunge un termine di ``inerzia'':
\[
  y_{k+1} = x_k - \frac{1}{L}\nabla f(x_k), \quad x_{k+1} = y_{k+1} + \frac{k-1}{k+2}(y_{k+1} - y_k).
\]
Converge come $O(L/k^2)$ invece di $O(L/k)$: \emph{optimal} per la classe dei metodi del primo ordine. Il termine di momento riutilizza la direzione dell'iterazione precedente.

\subsection{Line Search adattiva}

\subsubsection{Condizioni di Armijo e Wolfe}

\begin{defn}[Condizione di Armijo (sufficient decrease)]
Con direzione di discesa $d_k$ e step $\alpha > 0$:
\[
  f(x_k + \alpha d_k) \leq f(x_k) + c_1 \alpha \nabla f(x_k)^\top d_k, \quad c_1 \in (0, 1).
\]
\end{defn}

\begin{defn}[Condizione di curvatura (Wolfe)]
\[
  \nabla f(x_k + \alpha d_k)^\top d_k \geq c_2 \nabla f(x_k)^\top d_k, \quad c_2 \in (c_1, 1).
\]
\end{defn}

\begin{algobox}[Backtracking Armijo]
\begin{enumerate}
  \item $\alpha \leftarrow \alpha_0$ (es.\ $\alpha_0 = 1$), $\rho \in (0,1)$ (es.\ $\rho = 0.5$).
  \item Finché Armijo non è soddisfatta: $\alpha \leftarrow \rho\alpha$.
  \item Usa $\alpha$ per il passo $x_{k+1} = x_k + \alpha d_k$.
\end{enumerate}
\end{algobox}

\begin{thm}[Convergenza GD con Armijo]
Se $f$ è $L$-smooth e inferiormente limitata, e si usa backtracking Armijo, allora $\|\nabla f(x_k)\| \to 0$.
\end{thm}

\begin{remark}[Line search esatta per quadratiche]
Per $f(x) = \frac{1}{2}x^\top Q x - c^\top x$, $g_k = Qx_k - c$:
\[
  \alpha_k^* = \argmin_{\alpha \geq 0} f(x_k - \alpha g_k) = \frac{g_k^\top g_k}{g_k^\top Q g_k}.
\]
\end{remark}

%------------------------------------------------------------
\section{Metodo di Newton multivariato}

\begin{algobox}[Metodo di Newton]
Dato $x_0$, per $k = 0, 1, \ldots$:
\[
  H_f(x_k) d_k = -\nabla f(x_k), \quad x_{k+1} = x_k + d_k.
\]
(con eventuale line search: $x_{k+1} = x_k + \alpha_k d_k$)
\end{algobox}

\begin{intuition}
Newton \emph{minimizza il modello quadratico locale} di $f$:
\[
  m_k(d) = f(x_k) + \nabla f(x_k)^\top d + \tfrac{1}{2} d^\top H_f(x_k) d.
\]
Il minimo di $m_k$ (se $H_f \succ 0$) è esattamente $d_k = -H_f^{-1}\nabla f$.
\end{intuition}

\begin{thm}[Convergenza quadratica locale]
Se $f \in C^3$, $x^*$ è un minimo con $H_f(x^*) \succ 0$, e $x_0$ è sufficientemente vicino a $x^*$:
\[
  \|x_{k+1} - x^*\| \leq C \|x_k - x^*\|^2
\]
per una costante $C > 0$ che dipende da $\|\nabla^3 f\|$ e da $\lambda_{\min}(H_f(x^*))$.
\end{thm}

\textbf{Limitazioni:} Costo $O(n^3)$ per passo (fattorizzazione di $H_f$); $H_f$ potrebbe non essere $\succ 0$ lontano da $x^*$; non globalmente convergente senza modifiche.

\textbf{Modifiche pratiche:}
\begin{itemize}
  \item \textbf{Damped Newton}: $x_{k+1} = x_k + \alpha_k d_k$ con line search su $\alpha_k$.
  \item \textbf{Regularized Newton}: $H_f(x_k) + \mu_k I$ con $\mu_k > 0$ (connessione con trust-region).
  \item \textbf{Inexact Newton}: risolvere il sistema $H_f d = -\nabla f$ approssimativamente (es.\ con CG).
\end{itemize}

%------------------------------------------------------------
\section{Metodi quasi-Newton}

Idea: approssimare $H_f^{-1}$ tramite aggiornamenti di rango basso, usando solo informazioni sul gradiente.

\begin{defn}[Condizione secante]
Dati $s_k = x_{k+1} - x_k$ e $y_k = \nabla f(x_{k+1}) - \nabla f(x_k)$, si vuole $B_{k+1} s_k = y_k$ (approssimazione $B_{k+1} \approx H_f$).
\end{defn}

\subsection{Aggiornamento BFGS}

L'aggiornamento \textbf{BFGS} (Broyden--Fletcher--Goldfarb--Shanno), con $H_k \approx H_f^{-1}(x_k)$:
\begin{equation}
  H_{k+1} = \left(I - \rho_k s_k y_k^\top\right) H_k \left(I - \rho_k y_k s_k^\top\right) + \rho_k s_k s_k^\top,
  \quad \rho_k = \frac{1}{y_k^\top s_k}.
\end{equation}

\begin{prop}[Proprietà BFGS]
\begin{enumerate}
  \item Se $H_k \succ 0$ e $y_k^\top s_k > 0$ (garantito da Wolfe), allora $H_{k+1} \succ 0$.
  \item $H_{k+1} s_k = s_k / (y_k^\top s_k) \cdot (y_k^\top s_k) = s_k$... ovvero $H_{k+1} y_k = s_k$ (condizione secante).
  \item BFGS è l'aggiornamento di minima perturbazione (in norma di Frobenius pesata) che soddisfa la condizione secante.
\end{enumerate}
\end{prop}

\begin{thm}[Convergenza superlineare di BFGS]
Sotto ipotesi di regolarità standard (f $L$-smooth, $\tau$-strongly convex), BFGS ha convergenza \emph{superlineare}:
\[
  \|x_{k+1} - x^*\| = o(\|x_k - x^*\|).
\]
\end{thm}

\begin{remark}[\textbf{L-BFGS} per grande scala]
L-BFGS (Limited-memory BFGS) memorizza solo le ultime $m$ coppie $(s_k, y_k)$ per ricostruire $H_k^{-1}$ implicitamente. Costo $O(mn)$ per passo (anziché $O(n^2)$ per BFGS pieno). Fondamentale per ML: addestramento di reti neurali, ottimizzazione su $n = 10^7$ parametri.
\end{remark}

%------------------------------------------------------------
\section{Gradiente Coniugato per funzioni quadratiche}

Il metodo del \emph{gradiente coniugato} (CG) minimizza $f(x) = \frac{1}{2}x^\top Qx - c^\top x$ con $Q \succ 0$, equivalente a risolvere $Qx = c$.

\begin{defn}[$Q$-coniugazione]
Vettori $d_0, d_1, \ldots$ sono \emph{$Q$-coniugati} se $d_i^\top Q d_j = 0$ per $i \neq j$.
\end{defn}

\begin{algobox}[Gradiente Coniugato (CG)]
$x_0$ dato, $r_0 = c - Qx_0$, $d_0 = r_0$. Per $k = 0, 1, \ldots$:
\begin{align*}
  \alpha_k &= \frac{r_k^\top r_k}{d_k^\top Q d_k}, \quad &
  x_{k+1} &= x_k + \alpha_k d_k, \\
  r_{k+1} &= r_k - \alpha_k Q d_k, \quad &
  \beta_{k+1} &= \frac{r_{k+1}^\top r_{k+1}}{r_k^\top r_k}, \\
  d_{k+1} &= r_{k+1} + \beta_{k+1} d_k.&&
\end{align*}
\end{algobox}

\begin{thm}[Ottimalità di CG]
Dopo $k$ passi, $x_k$ minimizza $f$ sullo spazio di Krylov
\[
  \mathcal{K}_k(Q, r_0) = \mathrm{span}(r_0, Qr_0, \ldots, Q^{k-1}r_0).
\]
CG converge in al più $n$ passi (terminazione esatta).
\end{thm}

\begin{thm}[Convergenza lineare di CG]
\[
  \frac{\|x_k - x^*\|_Q}{\|x_0 - x^*\|_Q} \leq 2 \left(\frac{\sqrt{\kappa(Q)} - 1}{\sqrt{\kappa(Q)} + 1}\right)^k.
\]
CG è molto più veloce del gradient descent: il fattore è $\sqrt{\kappa}$ anziché $\kappa$.
\end{thm}

\begin{intuition}
CG minimizza il residuo relativo su uno spazio polinomiale: trova il polinomio $p_k$ di grado $\leq k$ con $p_k(0) = 1$ che minimizza $\max_{\lambda_i} |p_k(\lambda_i)|$. Se gli autovalori di $Q$ sono raggruppati in $m$ cluster, CG converge in $m$ passi.
\end{intuition}

%------------------------------------------------------------
\chapter{Ottimizzazione Non Vincolata — Metodi non-smooth}
%------------------------------------------------------------

\section{Funzioni non differenziabili}

Molte funzioni nel ML non sono differenziabili ovunque:
\begin{itemize}
  \item $f(x) = \|x\|_1 = \sum_i |x_i|$ (LASSO, sparsità, compressive sensing).
  \item $f(x) = \max_j g_j(x)$ (SVM, ottimizzazione worst-case).
  \item $f(x) = \max(0, x)$ (ReLU, attivazione neurale).
\end{itemize}

\begin{defn}[Subgradiente]
$g \in \R^n$ è un \emph{subgradiente} di $f$ convessa in $x$ se:
\[
  f(y) \geq f(x) + g^\top (y - x) \quad \forall y.
\]
L'insieme dei subgradienti $\partial f(x)$ è detto \emph{subdifferenziale}.
\end{defn}

\begin{example}
Per $f(x) = |x|$: $\partial f(x) = \{+1\}$ se $x > 0$, $[-1,+1]$ se $x = 0$, $\{-1\}$ se $x < 0$.
Per $f(x) = \|x\|_1$: $[\partial f(x)]_i = \sign(x_i)$ se $x_i \neq 0$, e $\in [-1,+1]$ se $x_i = 0$.
\end{example}

\begin{thm}
$x^*$ è un minimo globale di $f$ convessa $\Leftrightarrow$ $0 \in \partial f(x^*)$.
\end{thm}

\section{Metodo del subgradiente}

\begin{algobox}[Subgradient Method]
Dato $x_0$, step sizes $\{\alpha_k\}$:
\[
  x_{k+1} = x_k - \alpha_k g_k, \quad g_k \in \partial f(x_k).
\]
\end{algobox}

\begin{warningbox}
A differenza del gradient descent, il subgradiente \textbf{non garantisce} $f(x_{k+1}) < f(x_k)$: l'obiettivo può aumentare da un'iterazione all'altra. Si usa il record del best so far: $f_k^{\mathrm{best}} = \min_{i \leq k} f(x_i)$.
\end{warningbox}

\begin{thm}[Convergenza del subgradiente]
Se $f$ convessa, Lipschitz con costante $G$:
\begin{itemize}
  \item Con $\alpha_k = c/\sqrt{k}$ (dimishing): $f_k^{\mathrm{best}} - f^* \leq \frac{G \|x_0 - x^*\|}{\sqrt{k}}$ (tasso $O(1/\sqrt{k})$).
  \item Con $\alpha_k = c$ costante: $\liminf f_k^{\mathrm{best}} \leq f^* + Gc/2$.
\end{itemize}
\end{thm}

\begin{remark}
Il tasso $O(1/\sqrt{k})$ è \emph{ottimale} per metodi del primo ordine su funzioni Lipschitz non smooth. Per funzioni smooth, si può fare meglio: $O(1/k)$ con GD.
\end{remark}

\section{Metodo del gradiente prossimale}

Molte funzioni hanno la forma $f(x) = h(x) + r(x)$, con $h$ smooth e $r$ non-smooth ma semplice (es.\ $r(x) = \lambda \|x\|_1$).

\begin{defn}[Operatore prossimale]
Per $r: \R^n \to \R \cup \{+\infty\}$ e $\alpha > 0$:
\[
  \prox_{\alpha r}(v) = \argmin_x \left(r(x) + \frac{1}{2\alpha}\|x - v\|^2\right).
\]
\end{defn}

\begin{example}[Prox dell'$\ell^1$: soft thresholding]
Per $r(x) = \lambda\|x\|_1$:
\[
  [\prox_{\alpha\lambda\|\cdot\|_1}(v)]_i = \sign(v_i)\max(|v_i| - \alpha\lambda, 0).
\]
Questa è la \emph{sogliatura morbida} (soft thresholding).
\end{example}

\begin{algobox}[Proximal Gradient (ISTA)]
Per $f = h + r$ con $h$ $L$-smooth, $r$ convessa:
\[
  x_{k+1} = \prox_{\alpha r}\left(x_k - \alpha \nabla h(x_k)\right), \quad \alpha = 1/L.
\]
\end{algobox}

\begin{thm}[Convergenza ISTA]
Con $\alpha = 1/L$: $f(x_k) - f^* \leq \frac{L\|x_0 - x^*\|^2}{2k}$ (tasso $O(1/k)$).
\end{thm}

\begin{remark}[FISTA — accelerazione di Nesterov]
FISTA (Fast ISTA) aggiunge un passo di momento come Nesterov e raggiunge $O(1/k^2)$, ottimale per la classe. È il metodo standard per LASSO di grande scala.
\end{remark}

\section{Metodi bundle}

I metodi bundle sono la generalizzazione dei metodi del gradiente coniugato per funzioni non-smooth.

\begin{defn}[Modello bundle (piecewise-linear inferior)]
Dati gli iterati precedenti $x^1, \ldots, x^i$ con subgradienti $g^j \in \partial f(x^j)$, il \emph{bundle} è:
\[
  f_B(x) = \max_{j \leq i} \left\{f(x^j) + (g^j)^\top(x - x^j)\right\} \leq f(x).
\]
$f_B$ è un modello \emph{inferiore} di $f$.
\end{defn}

\begin{algobox}[Bundle Method (schema generale)]
Dati $x_0$, parametro di prossimità $\mu > 0$:
\begin{enumerate}
  \item Risolvi: $d^* = \argmin_d \{f_B(x_k + d) + \frac{\mu}{2}\|d\|^2\}$.
  \item Se $\|d^*\| \leq \eps$: STOP (criterio di ottimalità).
  \item Valuta $f(x_k + d^*)$ e $g^* \in \partial f(x_k + d^*)$.
  \item \textbf{Serious step}: se $f(x_k + d^*) \leq f(x_k) - m_1(f_B(x_k + d^*) - f(x_k))$: $x_{k+1} \leftarrow x_k + d^*$.
  \item \textbf{Null step}: altrimenti $x_{k+1} \leftarrow x_k$. Aggiorna il bundle con il nuovo punto.
\end{enumerate}
\end{algobox}

\begin{thm}[Convergenza dei metodi bundle]
Per $f$ convessa Lipschitz, con $\mu$ fisso e bounded: se $f^* > -\infty$, allora $\|d^*_i\| \to 0$, il che implica la convergenza all'ottimo.
\end{thm}

%------------------------------------------------------------
\chapter{Ottimizzazione Vincolata}
%------------------------------------------------------------

\section{Formulazione e topologia}

\begin{defn}[Problema vincolato]
\[
  \min_{x \in \R^n} f(x) \quad \text{s.t.} \quad g_i(x) \leq 0, \; i = 1,\ldots,m; \quad h_j(x) = 0, \; j = 1,\ldots,p.
\]
\end{defn}

\begin{thm}[Weierstrass]
Se $f$ è continua e $S \subseteq \R^n$ è compatto e non vuoto, il minimo è attinto.
\end{thm}

\begin{remark}
Per il problema vincolato, un punto stazionario di $f$ su $\R^n$ potrebbe essere inammissibile. Le condizioni di ottimalità richiedono una nozione diversa.
\end{remark}

\section{Condizioni KKT}

\subsection{Vincoli attivi e LICQ}

\begin{defn}[Vincoli attivi]
All'ottimo $x^*$, il vincolo $g_i$ è \emph{attivo} se $g_i(x^*) = 0$, altrimenti \emph{inattivo} ($g_i(x^*) < 0$).
\end{defn}

\begin{defn}[LICQ — Linear Independence Constraint Qualification]
Condizione LICQ: i gradienti dei vincoli attivi in $x^*$ sono linearmente indipendenti.
\end{defn}

\begin{thm}[Condizioni KKT — Karush-Kuhn-Tucker]
Se $x^*$ è un minimo locale e vale LICQ, esistono moltiplicatori $\mu^* \in \R^m$, $\lambda^* \in \R^p$ tali che:
\begin{align}
  &\nabla f(x^*) + \sum_{i=1}^m \mu_i^* \nabla g_i(x^*) + \sum_{j=1}^p \lambda_j^* \nabla h_j(x^*) = 0 \tag{stazionarietà} \\
  &g_i(x^*) \leq 0, \quad h_j(x^*) = 0 \tag{ammissibilità} \\
  &\mu_i^* \geq 0 \tag{non negatività} \\
  &\mu_i^* g_i(x^*) = 0 \quad \forall i \tag{complementarità}
\end{align}
\end{thm}

\begin{intuition}
Al minimo vincolato, il gradiente di $f$ deve trovarsi nel \emph{cono normale} all'insieme ammissibile in $x^*$: non si può migliorare $f$ lungo nessuna direzione ammissibile.

La condizione di complementarità $\mu_i^* g_i(x^*) = 0$ dice: o il vincolo $g_i$ è inattivo (e il suo moltiplicatore è zero), oppure il moltiplicatore è positivo (e il vincolo è attivo, in equilibrio).
\end{intuition}

\begin{thm}[Condizioni del secondo ordine]
Siano soddisfatte le KKT in $x^*$. Definiamo il \emph{Lagrangiano ridotto}:
\[
  L(x) = f(x) + \sum_i \mu_i^* g_i(x) + \sum_j \lambda_j^* h_j(x).
\]
\begin{itemize}
  \item \textbf{Necessaria}: $d^\top \nabla^2 L(x^*) d \geq 0$ per ogni $d$ nel cono critico.
  \item \textbf{Sufficiente}: $d^\top \nabla^2 L(x^*) d > 0$ per ogni $d \neq 0$ nel cono critico $\Rightarrow$ $x^*$ minimo locale isolato.
\end{itemize}
\end{thm}

\subsection{Lemma di Farkas e dualità}

\begin{lemma}[Farkas]
Per $A \in \R^{m \times n}$, $c \in \R^n$, esattamente una delle seguenti è vera:
\begin{enumerate}
  \item Esiste $x \geq 0$ con $Ax = c$.
  \item Esiste $y$ con $A^\top y \geq 0$ e $c^\top y < 0$.
\end{enumerate}
\end{lemma}

Il Lemma di Farkas è lo strumento fondamentale nella dimostrazione delle condizioni KKT.

%------------------------------------------------------------
\section{Dualità Lagrangiana}

\begin{defn}[Funzione Lagrangiana]
\[
  L(x, \mu, \lambda) = f(x) + \mu^\top g(x) + \lambda^\top h(x), \quad \mu \geq 0.
\]
\end{defn}

\begin{defn}[Funzione duale di Lagrange]
\[
  q(\mu, \lambda) = \inf_{x \in \R^n} L(x, \mu, \lambda).
\]
$q$ è sempre concava (anche se $f, g, h$ non lo sono).
\end{defn}

\begin{thm}[Weak duality]
Per ogni $x$ ammissibile e $(\mu \geq 0, \lambda)$: $q(\mu, \lambda) \leq f^*$.

Il \emph{problema duale} $\max_{\mu \geq 0, \lambda} q(\mu, \lambda)$ fornisce sempre un lower bound su $f^*$.
\end{thm}

\begin{defn}[Duality gap]
$f^* - q^* \geq 0$. Si ha \emph{strong duality} quando $f^* = q^*$.
\end{defn}

\begin{thm}[Condizione di Slater (strong duality per problemi convessi)]
Se il problema è convesso ($f, g_i$ convesse, $h_j$ affini) ed esiste $\hat{x}$ strettamente ammissibile ($g_i(\hat{x}) < 0$ per ogni $i$, $h_j(\hat{x}) = 0$ per ogni $j$), allora vale la strong duality: $f^* = q^*$.
\end{thm}

\begin{example}[Dualità LP]
$\min_x c^\top x$ s.t.\ $Ax = b$, $x \geq 0$ ha duale $\max_y b^\top y$ s.t.\ $A^\top y \leq c$. Per LP, la strong duality vale sempre (quando entrambi sono ammissibili).
\end{example}

\begin{example}[Dualità QP]
Per $\min \frac{1}{2}x^\top Q x + c^\top x$ s.t.\ $Ax = b$, $Q \succ 0$:
\[
  q(\lambda) = -\frac{1}{2}(c + A^\top\lambda)^\top Q^{-1}(c + A^\top\lambda) - b^\top\lambda.
\]
Funzione duale quadratica concava, strong duality vale.
\end{example}

%------------------------------------------------------------
\section{Algoritmi per l'ottimizzazione vincolata}

\subsection{Gradiente proiettato}

\begin{defn}[Proiezione su insieme convesso]
$\Pi_C(x) = \argmin_{y \in C} \|y - x\|^2$. Unica per $C$ convesso chiuso.
\end{defn}

\begin{algobox}[Projected Gradient Descent]
\[
  x_{k+1} = \Pi_S(x_k - \alpha_k \nabla f(x_k)).
\]
Converge come GD standard se $f$ è smooth e convessa, $S$ è convesso.
\end{algobox}

\begin{example}[Proiezioni facili]
\begin{itemize}
  \item Simplesso $\Delta = \{x \geq 0: \sum x_i = 1\}$: $O(n \log n)$.
  \item Box $[l, u]$: $\Pi_{[l,u]}(x)_i = \max(l_i, \min(u_i, x_i))$.
  \item Palla $\|x\| \leq r$: $\Pi(v) = r \cdot v / \max(r, \|v\|)$.
\end{itemize}
\end{example}

\subsection{Metodo di Frank-Wolfe (Conditional Gradient)}

\begin{algobox}[Frank-Wolfe]
\begin{enumerate}
  \item Linearizza $f$ in $x_k$: $s_k = \argmin_{x \in C} \nabla f(x_k)^\top x$.
  \item Aggiorna: $x_{k+1} = x_k + \alpha_k (s_k - x_k)$, $\alpha_k = 2/(k+2)$.
\end{enumerate}
Costo per passo: solo la minimizzazione lineare su $C$ (molto più facile della proiezione per certi insiemi, es.\ poliedri).
\end{algobox}

\begin{thm}[Convergenza Frank-Wolfe]
Per $f$ convessa $L$-smooth su $C$ compatto (diametro $D$):
\[
  f(x_k) - f^* \leq \frac{2LD^2}{k+2}.
\]
\end{thm}

\subsection{Metodi barriera (Interior Point)}

Idea: approssima il vincolo $g(x) \leq 0$ con la \emph{funzione barriera logaritmica} $-t\log(-g(x))$:
\[
  \min_x f(x) - t \sum_i \log(-g_i(x)) \quad \text{(nessun vincolo!)}.
\]
Si risolve una sequenza di problemi al diminuire di $t \to 0$. Il percorso delle soluzioni al variare di $t$ è detto \emph{central path}. I metodi interior point hanno complessità polinomiale e sono la scelta per LP/QP/SOCP di grande scala.

%============================================================
\part{Analisi Numerica Lineare (NLA)}

%------------------------------------------------------------
\chapter{Background: Algebra Lineare}
%------------------------------------------------------------

\section{Prodotto matriciale e interpretazioni}

\begin{remark}[Costo computazionale]
Prodotto $Av$ con $A \in \R^{m \times n}$: $O(mn)$ flop. Prodotto $AB$ con $B \in \R^{n \times p}$: $O(mnp)$. \textbf{Regola fondamentale}: $A(Bv)$ costa $O(mn + np)$, mentre $(AB)v$ costa $O(mnp + mp)$ — non formare mai esplicitamente $AB$ se si vuole solo $ABv$.
\end{remark}

\section{Sistemi triangolari}

Un sistema $Lx = b$ con $L$ triangolare inferiore si risolve con \emph{forward substitution} in $O(n^2)$; $Ux = b$ con $U$ triangolare superiore con \emph{back substitution}.

\begin{warningbox}
Non calcolare mai $A^{-1}$ esplicitamente per risolvere $Ax = b$: il costo è $O(n^3)$ (come LU) ma il risultato è meno accurato. Usare \texttt{A\textbackslash b} in Matlab o \texttt{np.linalg.solve} in Python.
\end{warningbox}

\section{Autovalori e forme quadratiche}

\begin{thm}[Teorema spettrale]
Ogni $A = A^\top \in \R^{n \times n}$ si decompone come $A = Q\Lambda Q^\top$ con $Q$ ortogonale, $\Lambda = \diag(\lambda_1, \ldots, \lambda_n)$, $\lambda_i \in \R$.
\end{thm}

\begin{defn}[Definitezza]
$Q \succ 0$ (definita positiva) $\Leftrightarrow$ $\lambda_i > 0$ per ogni $i$ $\Leftrightarrow$ $x^\top Qx > 0$ per ogni $x \neq 0$.
\end{defn}

\begin{prop}
Per ogni $A \in \R^{m \times n}$: $A^\top A \succeq 0$. Se $A$ ha rango massimo ($\rank A = n \leq m$): $A^\top A \succ 0$.
\end{prop}

%------------------------------------------------------------
\chapter{Norme e Matrici Ortogonali}
%------------------------------------------------------------

\section{Norme vettoriali}

\begin{defn}
La famiglia delle \emph{$p$-norme}: $\|x\|_p = \left(\sum_i |x_i|^p\right)^{1/p}$ per $p \geq 1$. La norma euclidea è $\|\cdot\|_2$; la norma $\|\cdot\|_\infty = \max_i |x_i|$; la norma $\|\cdot\|_1 = \sum_i |x_i|$.
\end{defn}

\section{Matrici ortogonali}

\begin{defn}
$U \in \R^{m \times m}$ è \emph{ortogonale} se $U^\top U = UU^\top = I$.
\end{defn}

\begin{prop}
Le matrici ortogonali preservano la norma euclidea: $\|Ux\|_2 = \|x\|_2$.
Conseguenza: i valori singolari di $UA$ e $AV$ (con $U, V$ ortogonali) coincidono con quelli di $A$.
\end{prop}

\section{Norme di matrici}

\begin{defn}[Norma indotta]
La norma di matrice indotta dalla norma vettoriale $\|\cdot\|$ è:
\[
  \|A\| = \max_{\|x\|=1} \|Ax\| = \max_{x \neq 0} \frac{\|Ax\|}{\|x\|}.
\]
\end{defn}

\begin{thm}[Norma spettrale]
La norma indotta dalla norma euclidea è la \emph{norma spettrale}:
\[
  \|A\|_2 = \sigma_1(A) = \max_i \sigma_i(A).
\]
Per matrici ortogonali: $\|Q_1 A Q_2\|_2 = \|A\|_2$ (invarianza unitaria).
\end{thm}

\begin{defn}[Norma di Frobenius]
$\|A\|_F = \sqrt{\sum_{i,j} A_{ij}^2} = \sqrt{\sum_i \sigma_i^2}$. Anche la norma di Frobenius è unitariamente invariante.
\end{defn}

%------------------------------------------------------------
\chapter{Decomposizione a Valori Singolari (SVD)}
%------------------------------------------------------------

\section{Definizione e proprietà}

\begin{thm}[SVD — Singular Value Decomposition]
Per ogni $A \in \R^{m \times n}$ (con $m \geq n$), esistono:
\begin{itemize}
  \item $U \in \R^{m \times m}$ ortogonale (vettori singolari sinistri).
  \item $V \in \R^{n \times n}$ ortogonale (vettori singolari destri).
  \item $\Sigma = \diag(\sigma_1, \ldots, \sigma_n) \in \R^{m \times n}$ con $\sigma_1 \geq \sigma_2 \geq \cdots \geq \sigma_n \geq 0$.
\end{itemize}
tali che $A = U \Sigma V^\top$.
\end{thm}

\begin{proof}[Prova per $A$ quadrata]
La matrice $A^\top A$ è simmetrica semidefinita positiva. Sia $A^\top A = V D V^\top$ la sua decomposizione spettrale con $D = \diag(\sigma_1^2, \ldots, \sigma_n^2)$, $\sigma_i \geq 0$. Poniamo $\Sigma = \sqrt{D}$ e $U_r = AV_r \Sigma_r^{-1}$ (vettori singolari sinistri per i valori singolari non nulli). Si verifica che $U_r^\top U_r = \Sigma_r^{-1} V_r^\top A^\top A V_r \Sigma_r^{-1} = I$, e completando $U_r$ a una base ortonormale si ottiene $U$.
\end{proof}

\begin{defn}[SVD thin/ridotta]
Per $A \in \R^{m \times n}$ con $m \geq n$, la \emph{thin SVD} è $A = \hat{U} \Sigma V^\top$ con $\hat{U} \in \R^{m \times n}$ (colonne ortonormali). Sufficiente per la maggior parte delle applicazioni.
\end{defn}

\begin{thm}[Sottospazi fondamentali]
Dalla SVD $A = U\Sigma V^\top$ con $\rank A = r$:
\begin{itemize}
  \item $\mathrm{Im}(A) = \mathrm{span}(u_1, \ldots, u_r)$ (prime $r$ colonne di $U$).
  \item $\ker(A) = \mathrm{span}(v_{r+1}, \ldots, v_n)$ (ultimi $n-r$ colonne di $V$).
  \item $\mathrm{Im}(A^\top) = \mathrm{span}(v_1, \ldots, v_r)$.
  \item $\ker(A^\top) = \mathrm{span}(u_{r+1}, \ldots, u_m)$.
\end{itemize}
\end{thm}

\begin{defn}[Pseudoinversa di Moore-Penrose]
$A^+ = V \Sigma^+ U^\top$, dove $\Sigma^+ = \diag(\sigma_1^{-1}, \ldots, \sigma_r^{-1}, 0, \ldots, 0)$.
\end{defn}

\section{Teorema di Eckart-Young (approssimazione di basso rango)}

\begin{thm}[Eckart-Young, 1936]
Sia $A = U\Sigma V^\top$ e sia $A_k = \sum_{i=1}^k \sigma_i u_i v_i^\top$ il \emph{troncamento al rango $k$}. Allora:
\[
  \|A - A_k\|_2 = \sigma_{k+1}, \quad \|A - A_k\|_F = \sqrt{\sigma_{k+1}^2 + \cdots + \sigma_n^2},
\]
e $A_k$ è il miglior approssimante di rango $\leq k$ di $A$ in entrambe le norme.
\end{thm}

\begin{intuition}
$A_k$ cattura le $k$ ``direzioni principali'' di $A$ (quelle con maggiore varianza/energia). Fondamentale per: PCA, compressione di immagini, sistemi di raccomandazione.
\end{intuition}

\begin{example}[Matrice dei voti studenti]
$A \in \R^{s \times t}$ con $A_{ij}$ = voto dello studente $i$ al test $j$. Se $A \approx A_1 = \sigma_1 u_1 v_1^\top$ (rango 1), allora $A_{ij} \approx \sigma_1 (u_1)_i (v_1)_j$: il voto è prodotto di una ``abilità'' dello studente per una ``difficoltà'' del test. Se $A \approx A_2$, ci sono due ``dimensioni'' di abilità.
\end{example}

%------------------------------------------------------------
\chapter{Il Problema dei Minimi Quadrati}
%------------------------------------------------------------

\section{Formulazione}

\begin{defn}[Least Squares]
Dato $A \in \R^{m \times n}$ ($m \geq n$), $b \in \R^m$:
\[
  \min_{x \in \R^n} \|Ax - b\|_2^2.
\]
Se $m > n$, il sistema $Ax = b$ è sovradeterminato (più equazioni che incognite) — in generale non ha soluzione esatta.
\end{defn}

\begin{thm}[Soluzione geometrica]
$x^*$ è soluzione LS $\Leftrightarrow$ il residuo $r^* = Ax^* - b$ è ortogonale a $\mathrm{Im}(A)$, cioè $A^\top(Ax^* - b) = 0$.
\end{thm}

\begin{intuition}
$Ax^*$ è la \emph{proiezione ortogonale} di $b$ su $\mathrm{Im}(A)$. Il punto più vicino a $b$ nell'immagine di $A$ si trova proiettando.
\end{intuition}

\section{Equazioni normali}

\begin{thm}[Equazioni normali]
$x^*$ è soluzione LS $\Leftrightarrow$ soddisfa le \emph{equazioni normali}:
\[
  A^\top A x^* = A^\top b.
\]
\end{thm}

Se $A$ ha \emph{rango massimo} ($\rank A = n$, cioè $A^\top A \succ 0$), la soluzione è unica: $x^* = (A^\top A)^{-1} A^\top b$.

\section{Soluzione via SVD}

Dalla thin SVD $A = \hat{U}\Sigma V^\top$, la soluzione di norma minima è:
\[
  x^* = A^+ b = V \Sigma^+ \hat{U}^\top b = \sum_{i=1}^r \frac{\hat{u}_i^\top b}{\sigma_i} v_i.
\]

\begin{remark}[Valori singolari piccoli e regolarizzazione]
Se $A$ ha valori singolari piccoli, la soluzione $x^*$ è amplificata enormemente. La \emph{regolarizzazione di Tikhonov} (ridge regression) modifica il problema:
\[
  \min_x \|Ax - b\|^2 + \delta\|x\|^2 \quad \Rightarrow \quad x_\delta = (A^\top A + \delta I)^{-1} A^\top b = V \, \mathrm{diag}\!\left(\frac{\sigma_i}{\sigma_i^2 + \delta}\right) \hat{U}^\top b.
\]
Per $\delta > 0$, i valori singolari piccoli vengono smorzati: $\sigma_i/(\sigma_i^2 + \delta) \approx 1/\delta$ se $\sigma_i \ll \sqrt{\delta}$.
\end{remark}

%------------------------------------------------------------
\chapter{Fattorizzazione QR}
%------------------------------------------------------------

\section{Definizione e teorema}

\begin{thm}[Fattorizzazione QR]
Per ogni $A \in \R^{m \times n}$ con $m \geq n$, esistono:
\begin{itemize}
  \item $Q \in \R^{m \times m}$ ortogonale.
  \item $R \in \R^{m \times n}$ triangolare superiore.
\end{itemize}
tali che $A = QR$. La \emph{thin QR} è $A = \hat{Q}R$ con $\hat{Q} \in \R^{m \times n}$ (colonne ortonormali) e $R \in \R^{n \times n}$ triangolare superiore.
\end{thm}

\section{Riflettori di Householder}

Il metodo numericamente stabile per calcolare la QR.

\begin{defn}[Riflettore di Householder]
Dato un vettore $v \in \R^m$, il riflettore di Householder è:
\[
  H = I - 2\frac{vv^\top}{\|v\|^2} \in \R^{m \times m}.
\]
$H$ è simmetrica e ortogonale ($H = H^\top$, $H^2 = I$). Riflette rispetto all'iperpiano ortogonale a $v$.
\end{defn}

\begin{prop}[Azzeramento di colonne]
Dato $x \in \R^m$, si sceglie $v = x - \|x\|e_1$ (o $v = x + \|x\|e_1$ per stabilità numerica, scegliendo il segno opposto a $x_1$). Allora:
\[
  Hx = \|x\|\, e_1 = (\|x\|, 0, \ldots, 0)^\top.
\]
Questo azzerava tutti gli elementi sotto il primo.
\end{prop}

\begin{algobox}[Fattorizzazione QR via Householder]
Per $j = 1, \ldots, n$:
\begin{enumerate}
  \item Considera $x = A_{j:m, j}$ (colonna corrente a partire dalla diagonale).
  \item Costruisci $v_j$ tale che $H_j x = \|x\| e_1$.
  \item Aggiorna: $A \leftarrow H_j A$ (applica il riflettore all'intera matrice).
\end{enumerate}
Alla fine: $H_n \cdots H_1 A = R$ (triangolare sup.), quindi $Q = H_1 \cdots H_n$.
\textbf{Costo}: $2mn^2 - \frac{2}{3}n^3$ flop.
\end{algobox}

\begin{thm}[Backward stability della QR]
La fattorizzazione QR calcolata con riflettori di Householder soddisfa:
\[
  \tilde{Q}\tilde{R} = A + \Delta A, \quad \frac{\|\Delta A\|_F}{\|A\|_F} = O(u),
\]
dove $u \approx 2.2 \times 10^{-16}$ è la precisione di macchina. La QR di Householder è \emph{backward stable}.
\end{thm}

\section{Soluzione dei minimi quadrati con QR}

Con la thin QR $A = \hat{Q}R$:
\[
  \|Ax - b\|^2 = \|\hat{Q}Rx - b\|^2 = \|Rx - \hat{Q}^\top b\|^2 + \|(I - \hat{Q}\hat{Q}^\top)b\|^2.
\]
Il secondo termine non dipende da $x$. Quindi si risolve il sistema triangolare:
\[
  Rx^* = \hat{Q}^\top b.
\]
\textbf{Costo totale}: QR in $O(mn^2)$, risoluzione $Rx = c$ in $O(n^2)$.

\begin{warningbox}
Le \emph{equazioni normali} $A^\top Ax = A^\top b$ risolte con Cholesky hanno costo simile, ma l'errore numerico è proporzionale a $\kappa(A)^2$ anziché $\kappa(A)$. Preferire QR per problemi mal condizionati.
\end{warningbox}

%------------------------------------------------------------
\chapter{Condizionamento e Stabilità}
%------------------------------------------------------------

\section{Numero di condizionamento}

\begin{defn}[Numero di condizionamento assoluto e relativo]
Per $f: \R^m \to \R^n$:
\[
  \kappa_{\mathrm{abs}}(f, x) = \lim_{\delta \to 0} \sup_{\|d\| \leq \delta} \frac{\|f(x+d) - f(x)\|}{\|d\|}, \quad
  \kappa_{\mathrm{rel}}(f, x) = \kappa_{\mathrm{abs}} \cdot \frac{\|x\|}{\|f(x)\|}.
\]
\end{defn}

\begin{thm}[Numero di condizionamento di un sistema lineare]
Per $f(b) = A^{-1}b$ (risolvere $Ax = b$):
\[
  \boxed{\kappa(A) = \|A\| \cdot \|A^{-1}\| = \frac{\sigma_1}{\sigma_n}}.
\]
Perturbando $b$ di $\delta b$: $\frac{\|\delta x\|}{\|x\|} \leq \kappa(A) \frac{\|\delta b\|}{\|b\|}$.
\end{thm}

\begin{intuition}
$\kappa(A) \gg 1$: problema \emph{mal condizionato}. Se $\kappa(A) = 10^6$ e $b$ ha un errore relativo di $10^{-12}$ (precisione di macchina), la soluzione $x$ ha un errore relativo di $\approx 10^{-6}$: si perdono 6 cifre decimali.
\end{intuition}

\begin{prop}[Distanza dalla singolarità]
$1/\kappa(A) = \min_{\tilde{A}\ \text{singolare}} \|\tilde{A} - A\|/\|A\|$.
Matrici con $\kappa(A) \gg 1$ sono ``quasi singolari''.
\end{prop}

\begin{thm}[Condizionamento del problema LS]
Per $\min_x \|Ax - b\|$ con $A$ a rango massimo e soluzione $x^*$:
\begin{itemize}
  \item Perturbazione di $b$: errore relativo proporzionale a $\kappa(A)$.
  \item Perturbazione di $A$: errore proporzionale a $\kappa(A)^2$ se $\|r^*\|/\|b\|$ non è trascurabile.
\end{itemize}
Nota: $\kappa(A^\top A) = \kappa(A)^2$, confermando l'instabilità delle equazioni normali.
\end{thm}

\section{Aritmetica floating point}

I numeri double-precision (IEEE 754, 64 bit) hanno struttura:
\[
  \underbrace{1 \text{ bit}}_{\text{segno}} \underbrace{11 \text{ bit}}_{\text{esponente}} \underbrace{52 \text{ bit}}_{\text{mantissa}}.
\]
\begin{itemize}
  \item Range: $\approx [10^{-308}, 10^{308}]$.
  \item Precisione di macchina: $u = 2^{-52} \approx 2.2 \times 10^{-16}$.
  \item Ogni numero reale $x$ è rappresentato da $\tilde{x}$ con $|\tilde{x} - x|/|x| \leq u$.
  \item Ogni operazione aritmetica introduce un errore relativo $\leq u$: $a \oplus b = (a+b)(1+\delta)$, $|\delta| \leq u$.
\end{itemize}

\begin{remark}[Problemi di cancellazione]
Se si sottraggono due numeri quasi uguali, le cifre significative vengono annullate. Es.\ $1.000001 - 1.000000 = 0.000001$: con 7 cifre significative, il risultato ne ha solo 1. Algoritmi che evitano la cancellazione sono preferibili.
\end{remark}

\section{Stabilità degli algoritmi}

\begin{defn}[Backward stability]
Un algoritmo che calcola $\tilde{y} \approx f(x)$ è \emph{backward stable} se:
\[
  \tilde{y} = f(\tilde{x}), \quad \frac{\|\tilde{x} - x\|}{\|x\|} = O(u).
\]
L'output è l'esatto risultato per un input \emph{leggermente perturbato}.
\end{defn}

\begin{thm}[Backward stability del prodotto interno]
Sia $y = a^\top b = \sum_{i=1}^n a_i b_i$ calcolato in macchina come $\tilde{y}$. Allora:
\[
  \tilde{y} = \tilde{a}^\top b = a^\top \tilde{b},
\]
dove $\|\tilde{a} - a\| \leq nu\|a\| + O(u^2)$ e analogamente per $b$. Il prodotto interno è backward stable.
\end{thm}

\begin{proof}[Prova (sketch)]
Ogni moltiplicazione $a_i \odot b_i = a_i b_i(1+\delta_i)$, $|\delta_i| \leq u$. Le addizioni successive introducono ulteriori errori. Un'analisi rigorosa mostra che l'errore accumulato soddisfa:
\[
  |\tilde{y} - y| \leq nu \sum_i |a_i b_i| + O(u^2) \leq nu \|a\|_1\|b\|_\infty.
\]
L'algoritmo può essere visto come il calcolo esatto di $\tilde{a}^\top b$ con $\tilde{a}_i = a_i(1+\varepsilon_i)$, $|\varepsilon_i| \leq (2i-1)u$: backward stable.
\end{proof}

\begin{example}[Prodotto interno di vettori lunghi]
Con $n = 1000$ e $u \approx 10^{-16}$: l'errore relativo del prodotto interno è $\leq 1000 \cdot 10^{-16} = 10^{-13}$. Ancora molto piccolo! Il prodotto interno è numericamente robusto.
\end{example}

\begin{thm}[Backward stability di Householder QR]
La fattorizzazione QR tramite riflettori di Householder è backward stable: $\tilde{Q}\tilde{R} = A + E$ con $\|E\| = O(u\|A\|)$.
\end{thm}

\begin{warningbox}
Il metodo di Gram-Schmidt classico \textbf{non} è backward stable: gli errori di cancellazione possono rendere le colonne di $Q$ non ortogonali. Gram-Schmidt modificato è meglio, ma Householder rimane il metodo preferito.
\end{warningbox}

\subsection{Test a posteriori: analisi del residuo}

\begin{thm}[Residual bound per sistemi lineari]
Dato $\tilde{x} \approx A^{-1}b$ con residuo $r = A\tilde{x} - b$:
\[
  \frac{\|\tilde{x} - x^*\|}{\|x^*\|} \leq \kappa(A) \cdot \frac{\|r\|}{\|b\|}.
\]
\end{thm}

Se $\|r\|/\|b\| = O(u)$ (residuo piccolo in machine precision), allora l'algoritmo è backward stable. Ma l'errore forward può ancora essere grande se $\kappa(A) \gg 1$.

%------------------------------------------------------------
\chapter{Metodi Iterativi per Sistemi Lineari}
%------------------------------------------------------------

\section{Motivazione: sistemi sparsi di grande scala}

Per $n = 10^5$--$10^7$: LU richiede $O(n^3) \gg 10^{15}$ flop — completamente impraticabile. Ma spesso la matrice è \emph{sparsa}: $\mathrm{nnz}(A) \ll n^2$ (es.\ $O(n)$ nonzeri per riga). Il prodotto $Av$ ha allora costo $O(\mathrm{nnz}(A))$.

I \emph{metodi iterativi di Krylov} sfruttano questa struttura: ogni passo richiede solo un prodotto matrice-vettore.

\section{Spazi di Krylov}

\begin{defn}
$\mathcal{K}_k(A, b) = \mathrm{span}(b, Ab, A^2b, \ldots, A^{k-1}b) \subseteq \R^n$.
\end{defn}

\begin{prop}
$x_k = q_{k-1}(A)b$ per qualche polinomio $q_{k-1}$ di grado $\leq k-1$. I metodi di Krylov scelgono $q_{k-1}$ in modo ottimale rispetto a qualche criterio.
\end{prop}

\section{Gradiente Coniugato per sistemi SPD}

CG (già visto nell'ottimizzazione) risolve $Qx = v$ con $Q \succ 0$:

\begin{thm}[Convergenza CG]
\[
  \frac{\|x_k - x^*\|_Q}{\|x_0 - x^*\|_Q} \leq 2\left(\frac{\sqrt{\kappa(Q)}-1}{\sqrt{\kappa(Q)}+1}\right)^k.
\]
Se $Q$ ha $m$ autovalori distinti, CG converge in esattamente $m$ passi. Per autovalori \emph{raggruppati} (clustered), la convergenza è rapida: il polinomio di Chebyshev approssima bene la ``forma'' dello spettro.
\end{thm}

\subsection{Precondizionamento}

\begin{defn}[Preconditioned CG]
$P \approx Q$ invertibile, $P \succ 0$. Si risolve $P^{-1}Qx = P^{-1}v$, che ha lo stesso spettro ``rimodellato''.
\end{defn}

La convergenza dipende da $\kappa(P^{-1}Q) \ll \kappa(Q)$. Scelte comuni:
\begin{itemize}
  \item $P = \diag(Q)$ (Jacobi): semplice, smorzamento delle scale.
  \item $P = $ fattore di Cholesky incompleto (ICC): più efficace.
  \item $P = $ ILUTP per sistemi non simmetrici.
\end{itemize}

\section{Algoritmo di Arnoldi}

Per sistemi non simmetrici, si costruisce una base ortonormale di $\mathcal{K}_k(A, b)$ stabilmente.

\begin{algobox}[Arnoldi]
$q_1 = b/\|b\|$. Per $j = 1, \ldots, k$:
\begin{enumerate}
  \item $w = Aq_j$.
  \item Per $i = 1, \ldots, j$: $h_{ij} = q_i^\top w$; $w \leftarrow w - h_{ij} q_i$ (ortogonalizzazione).
  \item $h_{j+1,j} = \|w\|$; se $h_{j+1,j} = 0$: STOP (breakdown). Altrimenti $q_{j+1} = w/h_{j+1,j}$.
\end{enumerate}
\end{algobox}

\begin{thm}[Fattorizzazione di Arnoldi]
Dopo $k$ passi: $AQ_k = Q_{k+1} H_k$, dove:
\begin{itemize}
  \item $Q_k = [q_1, \ldots, q_k] \in \R^{n \times k}$ ha colonne ortonormali.
  \item $H_k \in \R^{(k+1) \times k}$ è una matrice di Hessenberg superiore.
\end{itemize}
\end{thm}

\begin{intuition}
Arnoldi è Gram-Schmidt applicato alle potenze di $A$: costruisce una base ortogonale senza formare esplicitamente $[b, Ab, \ldots, A^{k-1}b]$ (che sarebbero numericamente instabili per potenze elevate).
\end{intuition}

\begin{remark}[Caso simmetrico: Lanczos]
Se $A = A^\top$, la matrice $H_k$ è simmetrica tridiagonale. L'algoritmo si semplifica nella ricorrenza di Lanczos: ogni passo richiede solo 2 prodotti matrice-vettore e la memorizzazione di 3 vettori (invece di $k$ per Arnoldi). Si ricade nel metodo CG per sistemi SPD.
\end{remark}

\section{GMRES}

\begin{defn}[GMRES — Generalized Minimal Residual]
Al passo $k$, GMRES trova la migliore approssimazione nello spazio di Krylov rispetto al residuo:
\[
  x_k = \argmin_{x \in x_0 + \mathcal{K}_k(A, r_0)} \|Ax - b\|.
\]
\end{defn}

\begin{algobox}[GMRES]
\begin{enumerate}
  \item Esegui $k$ passi di Arnoldi su $(A, r_0)$: ottieni $Q_{k+1}$, $H_k$.
  \item Nota che $r_k = Ax_k - b = Q_{k+1}(H_k z_k - \|r_0\|e_1)$ con $x_k = x_0 + Q_k z_k$.
  \item Risolvi il piccolo problema LS: $z_k = \argmin_z \|H_k z - \|r_0\|e_1\|$ (costo $O(k^2)$).
  \item Soluzione: $x_k = x_0 + Q_k z_k$.
\end{enumerate}
Il residuo si calcola come by-product dalla fattorizzazione QR di $H_k$ (aggiornata a ogni passo).
\end{algobox}

\begin{thm}[Convergenza GMRES]
Se $A = V\Lambda V^{-1}$ (diagonalizzabile), allora:
\[
  \frac{\|r_k\|}{\|r_0\|} \leq \kappa(V) \min_{p \in \mathcal{P}_k,\, p(0)=1} \max_{\lambda_i} |p(\lambda_i)|,
\]
dove $\mathcal{P}_k$ è l'insieme dei polinomi di grado $\leq k$. Il termine $\kappa(V)$ misura quanto $A$ ``non è normale''.
\end{thm}

\begin{intuition}
GMRES converge rapidamente se gli autovalori di $A$ sono \emph{raggruppati} in pochi cluster: il polinomio ottimale può essere piccolo su tutti gli autovalori con poche ``radici'' (una per cluster). Se $A$ ha autovalori ben separati, la convergenza è lenta.

\textbf{Breakdown fortunato}: se $\|w\| = 0$ in Arnoldi, allora $Ax = b$ ha una soluzione esatta nel Krylov — GMRES ha trovato la soluzione!
\end{intuition}

\begin{warningbox}
GMRES ha costo crescente per iterazione: $O(k^2)$ per la fattorizzazione QR di $H_k$ e $O(kn)$ per lo storage. Si usa \emph{restarted GMRES}: dopo $m$ passi si ricomincia da $x_m$. Il costo diventa $O(mn)$ per ciclo, ma la convergenza può rallentare.
\end{warningbox}

%------------------------------------------------------------
\chapter{Fattorizzazioni Dirette per Sistemi Lineari}
%------------------------------------------------------------

\section{Fattorizzazione LU}

\begin{thm}[LU con pivoting parziale]
Per ogni $A \in \R^{n \times n}$, con un'opportuna matrice di permutazione $P$:
\[
  PA = LU,
\]
con $L$ triangolare inferiore con diagonale unitaria ($|L_{ij}| \leq 1$), $U$ triangolare superiore. \textbf{Costo}: $\frac{2}{3}n^3$ flop.
\end{thm}

\begin{intuition}
LU è la fattorizzazione dell'eliminazione gaussiana. Il pivoting parziale (scegliere il pivot come l'elemento di modulo massimo nella colonna corrente) garantisce la stabilità numerica: $|L_{ij}| \leq 1$ previene la crescita degli errori di arrotondamento.
\end{intuition}

\begin{algobox}[Risoluzione via LU]
Per $Ax = b$ con $PA = LU$:
\begin{enumerate}
  \item Permuta: $\tilde{b} = Pb$.
  \item Forward substitution: $Ly = \tilde{b}$, $O(n^2)$.
  \item Back substitution: $Ux = y$, $O(n^2)$.
\end{enumerate}
Costo totale: $\frac{2}{3}n^3$ (fattorizzazione) + $2n^2$ (soluzione). Stessa fattorizzazione per più termini noti.
\end{algobox}

\begin{thm}[Backward stability di LU con pivoting]
LU con pivoting parziale è backward stable: $\tilde{L}\tilde{U} = PA + E$ con $\|E\| = O(u\|A\|)$ (tipicamente — casi patologici esistono ma sono rari in pratica).
\end{thm}

\section{Fattorizzazione $LDL^\top$ per matrici simmetriche}

Per $M = M^\top$, l'eliminazione simmetrica produce $M = LDL^\top$: costo $\frac{1}{3}n^3$ (metà di LU pieno sfruttando la simmetria). Il pivoting di Bunch-Kaufman usa blocchi $1 \times 1$ e $2 \times 2$ in $D$ per stabilità.

\section{Fattorizzazione di Cholesky}

\begin{thm}[Cholesky]
Se $M = M^\top \succ 0$, esiste un'unica $R$ triangolare superiore con diagonale positiva tale che $M = R^\top R$. \textbf{Costo}: $\frac{1}{3}n^3$ — metà di LU!
\end{thm}

\begin{prop}
$M \succ 0 \Leftrightarrow$ tutti i pivot dell'eliminazione di Cholesky sono positivi. Quindi Cholesky può servire come test di definitezza positiva (se fallisce, $M \not\succ 0$).
\end{prop}

\begin{algobox}[Risoluzione con Cholesky]
Per $Mx = b$ con $M \succ 0$:
\begin{enumerate}
  \item Calcola $R = \mathrm{chol}(M)$, $\frac{1}{3}n^3$.
  \item Forward: $R^\top y = b$, $O(n^2)$.
  \item Back: $Rx = y$, $O(n^2)$.
\end{enumerate}
\end{algobox}

\section{Confronto e scelta dell'algoritmo}

\begin{center}
\renewcommand{\arraystretch}{1.4}
\begin{tabular}{llll}
\hline
\textbf{Struttura di $A$} & \textbf{Algoritmo} & \textbf{Costo} & \textbf{Backward stable} \\
\hline
Quadrata generica & LU + pivoting parz. & $\frac{2}{3}n^3$ & Sì (praticamente) \\
Simmetrica & $LDL^\top$ (Bunch-K.) & $\frac{1}{3}n^3$ & Sì \\
SPD & Cholesky & $\frac{1}{3}n^3$ & Sì (no pivoting!) \\
LS, $m \geq n$ & QR (Householder) & $2mn^2 - \frac{2}{3}n^3$ & Sì \\
SPD sparsa & CG (precondizionato) & $O(k \cdot \mathrm{nnz})$ & — \\
Generale sparsa & GMRES (precond.) & $O(k \cdot \mathrm{nnz})$ & — \\
\hline
\end{tabular}
\end{center}

\begin{remark}[Regola pratica]
Per sistemi densi: scegliere in base alla struttura (Cholesky se SPD, LU altrimenti). Per sistemi sparsi di grande scala ($n > 10^4$): usare CG (se SPD) o GMRES (altrimenti) con un buon precondizionatore. Non usare mai \texttt{inv(A)}.
\end{remark}

\section{Esempi pratici}

\begin{example}[Riordino per ridurre il fill-in]
Per matrici sparse, LU/Cholesky può produrre fattori molto più densi (\emph{fill-in}). Riordinare con Cuthill-McKee (\texttt{symrcm} in Matlab) riduce l'ampiezza di banda.
\begin{lstlisting}[language=Matlab]
p = symrcm(A);
R = chol(A(p, p));          % Cholesky riordinato
x_perm = R \ (R' \ b(p));
x(p) = x_perm;              % permutazione inversa
\end{lstlisting}
\end{example}

\begin{example}[Tikhonov via SVD in Python]
\begin{lstlisting}[language=Python]
import numpy as np
U, s, Vt = np.linalg.svd(A, full_matrices=False)
delta = 1e-4
d = s / (s**2 + delta)      # filtro di Tikhonov
x = Vt.T @ (d * (U.T @ b))  # soluzione regolarizzata
\end{lstlisting}
\end{example}

%------------------------------------------------------------
\chapter{Riepilogo: Connessioni tra NLA e Ottimizzazione}
%------------------------------------------------------------

Le due parti del corso sono strettamente interconnesse:

\begin{enumerate}
  \item \textbf{CG è sia ottimizzazione che algebra lineare}: risolvere $Qx = v$ con $Q \succ 0$ equivale a minimizzare $\frac{1}{2}x^\top Qx - v^\top x$. Stessa ricorrenza, due interpretazioni.

  \item \textbf{Il passo di Newton richiede un sistema lineare}: $H_f(x_k)d_k = -\nabla f(x_k)$. Per sistemi grandi, Newton ``inesatto'' usa CG o GMRES per risolvere il sistema approssimativamente.

  \item \textbf{Le equazioni normali del LS sono SPD}: $A^\top Ax = A^\top b$, risolto con Cholesky $O(\frac{1}{3}n^3)$, ma QR è preferibile per stabilità ($\kappa^2$ vs $\kappa$).

  \item \textbf{SVD e PCA}: la decomposizione a valori singolari è lo strumento fondamentale per la compressione dei dati, il calcolo della pseudoinversa, e l'analisi della sensitività.

  \item \textbf{Sistemi KKT}: i punti stazionari di problemi con vincoli di uguaglianza soddisfano il sistema saddle-point:
  \[
    \begin{pmatrix} H & A^\top \\ A & 0 \end{pmatrix}
    \begin{pmatrix} x \\ \lambda \end{pmatrix}
    =
    \begin{pmatrix} -\nabla f \\ 0 \end{pmatrix}.
  \]
  Questo è un sistema lineare simmetrico ma indefinito (non SPD): richiedono $LDL^\top$ con pivoting.

  \item \textbf{Regolarizzazione di Tikhonov}: il problema di regressione $\|Ax - b\|^2 + \delta\|x\|^2$ è un problema LS regolarizzato. La soluzione si ottiene con SVD: $x_\delta = V \mathrm{diag}(\sigma_i/(\sigma_i^2+\delta)) U^\top b$. Il parametro $\delta$ bilancia fit dei dati e norma della soluzione (bias-variance tradeoff in ML).
\end{enumerate}

%============================================================
\appendix
\chapter{Appendice: Richiami e Tabelle di Riferimento}
%============================================================

\section{Notazione}

\begin{center}
\renewcommand{\arraystretch}{1.3}
\begin{tabular}{ll}
\hline
Simbolo & Significato \\
\hline
$A \in \R^{m \times n}$ & Matrice $m \times n$ reale \\
$A^\top$ & Trasposta \\
$A^{-1}$ & Inversa ($A$ quadrata non singolare) \\
$A^+$ & Pseudoinversa di Moore-Penrose \\
$\|v\|_2$ & Norma euclidea \\
$\|A\|_2$ & Norma spettrale $= \sigma_1(A)$ \\
$\|A\|_F$ & Norma di Frobenius $= \sqrt{\sum \sigma_i^2}$ \\
$\kappa(A)$ & Numero di condizionamento $= \sigma_1/\sigma_n$ \\
$\sigma_i(A)$ & $i$-esimo valore singolare \\
$\lambda_i(A)$ & $i$-esimo autovalore \\
$Q \succ 0$ & $Q$ definita positiva \\
$Q \succeq 0$ & $Q$ semidefinita positiva \\
$\mathcal{K}_k(A,b)$ & Spazio di Krylov di ordine $k$ \\
$u \approx 2.2 \times 10^{-16}$ & Precisione di macchina (double) \\
$\partial f(x)$ & Subdifferenziale di $f$ in $x$ \\
$\prox_{\alpha f}(v)$ & Operatore prossimale \\
\hline
\end{tabular}
\end{center}

\section{Costi computazionali principali}

\begin{center}
\renewcommand{\arraystretch}{1.3}
\begin{tabular}{lll}
\hline
\textbf{Operazione} & \textbf{Costo (flop)} & \textbf{Note} \\
\hline
$Av$, $A \in \R^{m \times n}$ & $2mn$ & matrice-vettore \\
$A^\top A$ & $mn^2$ & non fare se evitabile \\
Sostituzione triangolare $n \times n$ & $n^2$ & forw./back. subst. \\
LU (con pivoting) & $\frac{2}{3}n^3$ & \\
Cholesky & $\frac{1}{3}n^3$ & $A$ SPD \\
QR (thin, $m \geq n$) & $2mn^2 - \frac{2}{3}n^3$ & Householder \\
SVD (thin, $m \geq n$) & $O(mn^2)$ & \\
Eigendecomp. $n \times n$ & $O(n^3)$ & \\
CG, $k$ passi & $k \cdot (2\,\mathrm{nnz}(A) + O(n))$ & \\
GMRES, $k$ passi & $k \cdot (2\,\mathrm{nnz}(A) + O(kn))$ & storage cresce! \\
\hline
\end{tabular}
\end{center}

\section{Tassi di convergenza degli algoritmi di ottimizzazione}

\begin{center}
\renewcommand{\arraystretch}{1.5}
\begin{tabular}{p{3.5cm}p{4.5cm}p{4cm}}
\hline
\textbf{Metodo} & \textbf{Tasso} & \textbf{Ipotesi} \\
\hline
GD, passo fisso & $O(L/k)$ & $f$ convessa $L$-smooth \\
GD, passo fisso & $\left(1-\frac{1}{\kappa}\right)^k$ & $f$ str. convessa \\
GD, quadratica & $\left(\frac{\kappa-1}{\kappa+1}\right)^k$ & $f$ quadratica, $\alpha$ ottimo \\
Nesterov/FISTA & $O(L/k^2)$ & $f$ convessa $L$-smooth \\
CG & $2\left(\frac{\sqrt{\kappa}-1}{\sqrt{\kappa}+1}\right)^k$ & $Q \succ 0$ \\
Newton (locale) & $O(\|e_k\|^2)$ quadratica & $f \in C^3$, vicino a $x^*$ \\
BFGS & $o(\|e_k\|)$ superlineare & $f$ smooth str. conv. \\
Subgradiente & $O(G/\sqrt{k})$ & $f$ convessa $G$-Lipschitz \\
ISTA/prox grad. & $O(L/k)$ & $f = h + r$, $h$ $L$-smooth \\
\hline
\end{tabular}
\end{center}

\section{Regole pratiche per l'esame}

\begin{keybox}[Scelta dell'algoritmo di ottimizzazione]
\begin{itemize}
  \item $f$ smooth convessa, $n$ piccolo: GD o Newton.
  \item $f$ smooth convessa, $n$ grande: L-BFGS o CG (per quadratiche).
  \item $f = h + r$ separabile smooth+nonsmooth: ISTA/FISTA (proximal gradient).
  \item $f$ nonsmooth in generale: metodo del subgradiente o bundle.
  \item Problema vincolato convesso: gradiente proiettato, Frank-Wolfe, o interior point.
\end{itemize}
\end{keybox}

\begin{keybox}[Scelta del metodo per sistemi lineari $Ax = b$]
\begin{itemize}
  \item $A$ SPD, densa: Cholesky ($\frac{1}{3}n^3$).
  \item $A$ simmetrica indef., densa: $LDL^\top$.
  \item $A$ generica, densa: LU con pivoting.
  \item $A$ SPD sparsa: PCG (CG precondizionato).
  \item $A$ sparsa non simm.: GMRES precondizionato.
  \item LS $\|Ax - b\|$, $A$ densa: QR (preferire a equazioni normali).
  \item LS con dati noisy: regolarizzazione Tikhonov via SVD.
\end{itemize}
\end{keybox}

%============================================================
\backmatter

\bibliographystyle{plain}
\begin{thebibliography}{99}

\bibitem{trefbau}
L.N.\ Trefethen e D.\ Bau III, \emph{Numerical Linear Algebra}, SIAM, 1997.

\bibitem{demmel}
J.W.\ Demmel, \emph{Applied Numerical Linear Algebra}, SIAM, 1997.

\bibitem{frangio25}
A.\ Frangioni, \emph{Optimization \& Learning --- Lecture Notes 2025/26}, Università di Pisa, 2025.

\bibitem{bertsekas}
D.P.\ Bertsekas, \emph{Nonlinear Programming}, 3rd ed., Athena Scientific, 2016.

\bibitem{nocedal}
J.\ Nocedal e S.J.\ Wright, \emph{Numerical Optimization}, 2nd ed., Springer, 2006.

\bibitem{boyd}
S.\ Boyd e L.\ Vandenberghe, \emph{Convex Optimization}, Cambridge University Press, 2004.

\bibitem{strang}
G.\ Strang, \emph{Linear Algebra and Learning from Data}, Wellesley-Cambridge Press, 2019.

\bibitem{bubeck}
S.\ Bubeck, \emph{Convex Optimization: Algorithms and Complexity}, arXiv:1405.4980, 2015.

\end{thebibliography}

\end{document}
